	\chapter{Applications of Integration}
	Integration is a powerful tool for calculating areas, volumes, and many other quantities. However, before we use an integral to represent or compute such a quantity, there is a more fundamental issue to settle: is the function we intend to integrate actually integrable?
	\section{More about Area}
	\begin{definitionenv}
		A \textbf{circular disk} of radius $r$ is the set of all points inside or on the boundary of a circle of radius $r$.
	\end{definitionenv}
	\begin{theoremenv}
		A circular disk is measurable and its area $A(r)$ is given by the integral
		\[
		A(r) = 2\int_{-r}^r \sqrt{r^2-x^2}\,\mathrm{d}x.
		\]
	\end{theoremenv}
	\begin{definitionenv}
		We define the number $\pi$ to be the area of a unit disk, i.e., the radius of which is exactly $1$. Symbolically: 
		\[
		\pi := 2\int_{-1}^1 \sqrt{1-x^2}\,\mathrm{d}x.
		\]
	\end{definitionenv}
	\begin{corollaryenv}
		The area of a circular disk with radius $r$ is given by
		\[
		A(r)=\pi r^2.
		\] 
	\end{corollaryenv}
	\begin{theoremenv}
		Let $f$ be nonnegative and integrable on $[a,b]$ and let $S$ be its ordinate set. We define a new function $g$ denoted by
		\[
		g(x):=kf\left( \frac{x}{k}\right) \qquad \text{with}\ x\in[ka,kb],\, k>0
		\]
		and its ordinate set $kS$. Then the area of $kS$ is given by
		\[
		A(kS)=k^2A(S).
		\]
	\end{theoremenv}
	\begin{remarkenv}
		In fact, this new $g$'s shape is ``similar" to $f$. Thus, we call this definition a similarity transformation.
	\end{remarkenv}
	\section{Integration formulas for the sine and cosine}
	We omit deriving properties of sine and cosine in this lecture note since I'm lazy. However, with the following theorems and lemmas, we can derive the integral formulas for sine and cosine.
	\begin{theoremenv}
		For $0<x<\dfrac{1}{2}\pi$, we have
		\[
		0<\cos x< \frac{\sin x}{x} < \frac{1}{\cos x}.
		\]
	\end{theoremenv}
	\begin{lemmaenv}
		Here are several inequalities important in deriving integral formulas
		for sine and cosine.
		\begin{enumerate}
			\item For any $x\in\mathbb R$,
			\begin{align*}
				2\sin\frac{x}{2}\sum_{k=1}^{n}\cos(kx)=\sin\!\left(\Big(n+\frac12\Big)x\right)-\sin\frac{x}{2}.
			\end{align*}
			
			\item For any $\theta\in\mathbb R$,
			\begin{align*}
				2\sin\theta\sum_{k=1}^{n}\cos(2k\theta)=\sin(2n+1)\theta-\sin\theta
				\\
				2\sin\theta\sum_{k=0}^{n-1}\cos(2k\theta)=\sin(2n-1)\theta+\sin\theta.
			\end{align*}
			
			\item Suppose $\theta$ satisfies $0<2n\theta<\dfrac12\pi$, we have
			\begin{align*}
				\sin(2n+1)\theta-\sin\theta<\frac{\sin\theta}{\theta}\sin(2n\theta).
			\end{align*}
			
			\item Suppose $\theta$ satisfies $0<2n\theta<\dfrac12\pi$. First we have
			\begin{align*}
				\frac{\cos(n-1)\theta}{\cos(n\theta)}>\frac{\sin\theta}{\theta}
			\end{align*}
			and then we can deduce
			\begin{align*}
				\frac{\sin\theta}{\theta}\sin(2n\theta)<\sin(2n-1)\theta+\sin\theta.
			\end{align*}
		\end{enumerate}
	\end{lemmaenv}
	\begin{theoremenv}
		For any $0<a\le \dfrac12\pi$, we have
		\begin{align*}
			\frac{a}{n}\sum_{k=1}^{n}\cos\frac{ka}{n}<\sin a
			<
			\frac{a}{n}\sum_{k=0}^{n-1}\cos\frac{ka}{n}.
		\end{align*}
	\end{theoremenv}
	\begin{theoremenv}
		For every real $a$ we have
		\begin{align*}
			\int_0^a \cos x\,\mathrm{d}x = \sin a,\qquad
			\int_0^a \sin x\, \mathrm{d}x = 1-\cos a.
		\end{align*}
		More generally, for every real $a,b$, we have
		\begin{align*}
			\int_a^b \cos x\,\mathrm{d}x = \sin x\,\Big|^b_a\, ,\qquad
			\int_a^b \sin x\, \mathrm{d}x = (-\cos x)\, \Big|^b_a\, .
		\end{align*}
	\end{theoremenv}

	\section{Polar Coordinates}
	\begin{definitionenv}
		The polar coordinates of a point $P$ in the plane are an ordered pair $(r,\theta)$ where $r$ is the distance from $P$ to the origin $O$ and $\theta$ is the angle formed by the line segment $OP$ and the positive $x$-axis. Here $r\ge 0$ and $\theta\in \mathbb{R}$. They are related to the rectangular coordinates $(x,y)$ of $P$ by the equations
		\[
		x=r\cos \theta,\quad y=r\sin \theta.
		\]
		Here we call $r$ the \textbf{radial distance} of $P$ and $\theta$ a \textbf{polar angle}.
	\end{definitionenv}
	\begin{remarkenv}
		Note that we call $\theta$ ``a'' polar angle instead of ``the'' polar angle because $\theta+2n\pi$ also represents the same angle for any integer $n$. But the radial distance $r$ is uniquely determined. 
		
		If $r=0$, then by convention $\theta$ can be any real number.
	\end{remarkenv}
	\begin{definitionenv}
		A \textbf{curve} is said to be given in polar coordinates if there is a nonnegative function $f$ defined on an interval $I$ such that every point $P$ on the curve has polar coordinates $(r,\theta)$ satisfying
		\[
		r=f(\theta),\quad \theta\in I.
		\]
		And this equation is called the \textbf{polar equation} of the curve. Similarly, the \textbf{graph} of $f$ in polar coordinates is the set of all points $P$ with polar coordinates $(r,\theta)$ satisfying the above equation.
	\end{definitionenv}

	\begin{remarkenv}
		In some textbooks, the variable $r$ in polar coordinates is allowed to be negative: when $r<0$, the point $(r,\theta)$ is identified with the point $(-r,\theta+\pi)$.  Note that the points $(r,\theta)$ and $(-r,\theta)$ are symmetric with respect to the origin. The additional definition allows the existence of inner loops in the graph of a polar equation. But in this lecture note, we only consider nonnegative functions in polar coordinates.
	\end{remarkenv}
	\begin{exampleenv}
		Sketch the curves for
		\[
		f(\theta)=1+c\sin\theta\quad \text{where}\quad\theta\in[0,2\pi],\ c\in\{-2,2\}
		\]
		under two different defintion and compare them.
	\end{exampleenv}
	\begin{definitionenv}
		Let $f$ be a nonnegative function defined on an interval $[a,b]$, where $0\le b-a\le 2\pi$. The set of all points $P$ in the plane with polar coordinates $(r,\theta)$ satisfying
		\[
		0\le r \le f(\theta),\quad \theta\in [a,b]
		\]
		is called the \textbf{radial set} determined by $f$.
	\end{definitionenv}
	\begin{definitionenv}
		Similarly as in rectangular coordinates, a nonnegative function $s$ defined on $[a,b]$ is called a \textbf{step function in polar coordinates} if there exists a partition $P=\left\{\theta_0,\theta_1,\cdots,\theta_n\right\}$ of $[a,b]$ such that
		\[
		s(\theta)=s_k\quad\text{where}\quad \theta \in (\theta_{k-1},\theta_k),\quad k=1,2,\cdots,n.
		\]
		And the area of this sector is given by
		\[
		A=\frac12 s_k^2(\theta_k-\theta_{k-1}).
		\]
		Since $b-a<2\pi$, there is no overlapping between different sectors. The area of the radial set determined by $s$ is given by
		\[
		A(S)=\frac12 \sum_{k=1}^n s_k^2(\theta_k-\theta_{k-1})=\frac12 \int_a^b s^2(\theta)\,\mathrm{d}\theta.
		\]
	\end{definitionenv}
	\begin{theoremenv}
		Let $f$ be a nonnegative function on $[a,b]$, where $0\le b-a\le 2\pi$. Assume the radial set $S$ determined by $f$ is measurable. If $f^2$ is integrable, then the area $A(S)$ is given by the integral
		\[
		A(S)=\frac12 \int_a^b f^2(\theta)\,\mathrm{d}\theta.
		\]
	\end{theoremenv}

	\section{Volume}
	Similarly, we can use axiomatic definitions to define volume like we did for area.

	\begin{definitionenv}
		We assume there exists a class $\mathcal{A}$ of solids and a set function $V$, whose domain is $\mathcal{A}$, with the following properties:
		\begin{enumerate}
			\item Nonnegative property: For each set $S$ in $\mathcal{A}$ we have $V(S) \ge 0$.
			
			\item Additive property: If $S$ and $T$ are in $\mathcal{A}$, then 
			$S \cup T$ and $S \cap T$ are in $\mathcal{A}$, and we have
			\[
				V(S \cup T) = V(S) + V(T) - V(S \cap T).
			\]
			
			\item Difference property: If $S$ and $T$ are in $\mathcal{A}$ with $S \subseteq T$, 
			then $T - S$ is in $\mathcal{A}$, and we have 
			\[
				V(T - S) = V(T) - V(S).
			\]
			
			\item Cavalieri's principle.: If $S$ and $T$ are two Cavalieri solids in $\mathcal{A}$ 
			with $A(S \cap F) \le A(T \cap F)$ for every plane $F$ perpendicular to a given line, 
			then $V(S) \le V(T)$.
			
			\item Choice of scale.: Every box $B$ is in $\mathcal{A}$. 
			If the edges of $B$ have lengths $a, b,$ and $c$, then 
			\[
				V(B) = abc.
			\]
			
			\item Every convex set is in $\mathcal{A}$.
		\end{enumerate}
	\end{definitionenv}
	\begin{definitionenv}[Additional]
		Some definitions are lost above and here we add it.
		\begin{enumerate}
			\item A solid $S$ in $\mathcal{A}$ is called a \textbf{Cavalieri solid} if for every plane $F$ perpendicular to a given line $L$, the area $A(S\cap F)$ of the cross section of $S$ by $F$ is well-defined, i.e., the cross section is measurable.
			\item A \textbf{box} or \textbf{rectangular parallelepiped} is any set congruent to a set of the form
			\[
			\{(x,y,z)\in \mathbb{R}^3\ |\ 0\le x \le a,\ 0\le y \le b,\ 0\le z \le c \},
			\]
			where nonnegative real numbers $a,b,c$ are called the lengths of the edges of the box.
			\item A solid $S$ in $\mathcal{A}$ is called a \textbf{convex set} if for every two points $P,Q\in S$, the line segment $PQ$ is also contained in $S$.
			\item A plane set is call \textbf{bounded} if it is a subset of some square in the plane.
		\end{enumerate}
	\end{definitionenv}
	\begin{corollaryenv}
		Here's some propositions for some basic cases.
		\begin{enumerate}
			\item The empty set $\varnothing$ is in $\mathcal{A}$ and $V(\varnothing)=0$.
			\item Monotonicity property: If $S$ and $T$ are in $\mathcal{A}$ with $S\subseteq T$, then $V(S)\le V(T)$.
			\item Every bounded plane set $S$ in $\mathcal{A}$ has volume zero, i.e., $V(S)=0$.
		\end{enumerate}
	\end{corollaryenv}
	\begin{remarkenv}
		Note that the Cavalieri's principle can be applied twice: if $A(S\cap F)=A(T\cap F)$ for every plane $F$ perpendicular to a given line, then $V(S)=V(T)$. In Chinese, we call this property Zu Geng's principle.
	\end{remarkenv}
	\begin{definitionenv}
		A \textbf{right cylindrical solid} is a set congruent to a set $S$ of the form
		\[
		S=\left\{(x,y,z)\in \mathbb{R}^3\ |\ (x,y)\in B,\, a\le z \le b \right\},
		\]
		where $B$ is a bounded plane set and $a,b$ are real numbers with $a\le b$. The base of the right cylindrical solid is the set $B$, and its height is $b-a$.
	\end{definitionenv}
	\begin{theoremenv}
		Let $S$ be a right cylindrical solid with base $B$ and height $h$. Then $S$ is in $\mathcal{A}$ and its volume is given by
		\[
		V(S)=A(B)\cdot h=\int_a^b A_S(z)\,\mathrm{d}z
		\]
		More generally, let $R$ be a Cavalieri solid in $\mathcal{A}$ with a cross-sectional area function $A_R$ which is integrable on an interval $[a,b]$ and zero outside $[a,b]$. Then the volume of $R$ is given by
		\[
		V(R)=\int_a^b A_R(z)\,\mathrm{d}z.
		\]
	\end{theoremenv}
	
	\begin{theoremenv}
	Let $f$ and $g$ be nonnegative integrable functions on $[a,b]$ satisfying $g(x) \le f(x)$ for all $x \in [a,b]$. If the region bounded by $y = f(x)$ and $y = g(x)$ is revolved about the $x$-axis, then the resulting solid of revolution (if it belongs to $\mathcal{A}$) has volume
	\[
		V = \pi \int_a^b \bigl[f^2(x) - g^2(x)\bigr] \, \mathrm{d}x .
	\]
	In particular, if $g(x) = 0$, the volume is
	\[
		V = \pi \int_a^b f^2(x) \, \mathrm{d}x .
	\]
	\end{theoremenv}
	\begin{exampleenv}
		 Let $f(x) = \sqrt{r^2 - x^2}$ on $[-r,r]$, then the solid is a sphere of radius $r$, and its volume is given by
		\[
		V = \pi \int_{-r}^{r} (r^2 - x^2)\, \mathrm{d}x = \frac{4}{3}\pi r^3.
		\]
	\end{exampleenv}

	\section{Indefinite Integral}
	\begin{definitionenv}
		Let $f$ be a real-valued function integrable on $[a,b]$. The \textbf{average value} of $f$, denoted by $\operatorname{avg}(f)$, on $[a,b]$ is defined to be
		\[
		\operatorname{avg}(f)=\frac{1}{b-a}\int_a^b f(x)\,\mathrm{d}x.
		\]
	\end{definitionenv}
    \begin{remarkenv}
		Note that the definition of average value of a function is consistent with the definition of average value of a finite set of numbers, i.e., the discrete average. You might not see how this definition works in this section, but it will be useful in later section.
	\end{remarkenv}
    \begin{theoremenv}[Weighted Mean Value Theorem for Integrals]
        Assume $f$ and $g$ are continuous on $[a,b]$. If $g$ never changes sign on $[a,b]$, then for some $c\in[a,b]$, we have
        \[
        \int_a^b f(x)g(x)\,\mathrm{d}x = f(c)\int_a^b g(x)\,\mathrm{d}x.
        \]
    \end{theoremenv}
    \begin{corollaryenv}[Mean Value Theorem for Integrals]
        Let $g(x)\equiv 1$, we have the \textbf{Mean Value Theorem for Integrals}:
        \[
        \int_a^b f(x)\,\mathrm{d}x = f(c)(b-a).
        \]
    \end{corollaryenv}
    \begin{remarkenv}
        The proof is in Chapter 4.
    \end{remarkenv}
	\begin{definitionenv}
		Let $f$ be a real-valued function integrable on every closed interval contained in an open interval $I$. An \textbf{indefinite integral} of $f$ on $I$ is the function $F$ defined by
		\[
		F(x)=\int_a^x f(t)\,\mathrm{d}t,\quad x\in I,
		\]
		where $a$ is an arbitrary point in $I$.
	\end{definitionenv}
	\begin{remarkenv}
		Indeed, a different choice of $a$ will only lead to a different indefinite integral by a constant, since
		\[
		\int_b^x f(t)\,\mathrm{d}t = \int_a^x f(t)\,\mathrm{d}t - \int_a^b f(t)\,\mathrm{d}t.
		\]
	\end{remarkenv}
	\begin{theoremenv}
		Suppose $f$ integrable and nonnegative on $[a,b]$, then the indefinite integral $A(f)$ is increasing.
	\end{theoremenv}
	\begin{definitionenv}
		A function $g$ is said to be \textbf{convex} on an interval $[a,b]$ if for every $x_1,x_2\in [a,b]$ and $0\le \lambda \le 1$, we have
		\[
		g\left(\lambda x_1 + (1-\lambda)x_2\right) \le \lambda g(x_1) + (1-\lambda)g(x_2).
		\]
		Similarly, $g$ is said to be \textbf{concave} on $[a,b]$ if the above inequality is reversed, i.e.,
		\[
		g\left(\lambda x_1 + (1-\lambda)x_2\right) \ge \lambda g(x_1) + (1-\lambda)g(x_2).
		\]
	\end{definitionenv}
	\begin{remarkenv}
		Geometrically, a function $g$ is convex on $[a,b]$ if the line segment joining any two points on the graph of $g$ lies above or on the graph between these two points. Similarly, $g$ is concave on $[a,b]$ if the line segment joining any two points on the graph of $g$ lies below or on the graph between these two points.
	\end{remarkenv}
	\begin{theoremenv}
		Let $f$ be integrable on $[a,b]$ and let $F$ be an indefinite integral of $f$ on $[a,b]$. Then:
        \begin{enumerate}
            \item if $f$ is increasing on $[a,b]$, then $F$ is convex on $[a,b]$;
            \item if $f$ is decreasing on $[a,b]$, then $F$ is concave on $[a,b]$.
        \end{enumerate}
	\end{theoremenv}

	\newpage
